\documentclass[11pt,a4paper]{article}

% --- packages ---
\usepackage[utf8]{inputenc}
\usepackage[T1]{fontenc}
\usepackage{geometry}
\geometry{margin=2.5cm}
\usepackage{graphicx}
\usepackage{booktabs}
\usepackage{hyperref}
\usepackage{amsmath}
\usepackage{amssymb}
\usepackage{float}
\usepackage{caption}
\usepackage{subcaption}
%\usepackage{enumitem}
\usepackage{xcolor}

\hypersetup{colorlinks=true, linkcolor=blue, citecolor=blue, urlcolor=blue}

\title{
  \textbf{SEDS 537 -- Machine Learning}\\[4pt]
  \Large Final Project Report\\[8pt]
  \large Multilingual Toxic Comment Classification\\
  with Parameter-Efficient Adapter Fine-Tuning
}
\author{
  Yasin T\"{u}rk -- 323011020\\
  Instructor: Prof.~Dr.~Aytu\u{g} Onan\\
  \.{I}zmir Y\"{u}ksek Teknoloji Enstit\"{u}s\"{u} -- Spring 2026
}
\date{}

\begin{document}
\maketitle

% =========================================================================
\section{Introduction}
% =========================================================================

Online platforms receive millions of comments every day in many different
languages. Detecting toxic comments automatically is important because manual
moderation cannot keep up with this volume. Most existing toxicity classifiers
work well only in English, so there is a clear need for multilingual solutions.

This project builds a \textbf{multilingual toxic comment classifier}. The main
idea is to use \textbf{mBERT} (multilingual BERT) with \textbf{adapter modules}
instead of updating all model parameters. Adapters are small bottleneck layers
inserted into the transformer blocks. During training, only these small layers
are updated while the rest of the model stays frozen. This approach is called
\textit{parameter-efficient fine-tuning}.

The key question we try to answer is: \textit{Can adapter-based fine-tuning
achieve similar cross-lingual performance to full fine-tuning while training
much fewer parameters?}

We compare six models in total: three classical machine learning baselines
(TF-IDF with Logistic Regression, Naive Bayes, and SVM), a BiLSTM neural
network, mBERT with full fine-tuning, and mBERT with adapters (our proposed
method).

% =========================================================================
\section{Dataset}
% =========================================================================

We use the \textbf{Jigsaw Multilingual Toxic Comment Classification} dataset
from Kaggle \cite{jigsaw}. The dataset has the following structure:

\begin{itemize}  \item \textbf{Training set:} 223,549 English comments with binary toxicity
        labels (toxic = 1, non-toxic = 0).
  \item \textbf{Validation set:} 8,000 multilingual comments in Turkish (3,000),
        Spanish (2,500), and Italian (2,500) with binary labels.
  \item \textbf{Test set:} 63,812 multilingual comments (used for the Kaggle
        competition; no public labels).
\end{itemize}

The training data has a strong \textbf{class imbalance}: only 9.6\% of comments
are toxic and 90.4\% are non-toxic. This imbalance makes it harder to detect
toxic comments because the models tend to predict everything as non-toxic.

We split the English training data into 80\% for training and 20\% for testing
(stratified by label). The multilingual validation set is used for
hyperparameter tuning, threshold tuning, and cross-lingual evaluation.

\subsection{Text Cleaning}

We use two different cleaning functions depending on the model:

\begin{itemize}  \item \textbf{Aggressive cleaning} (for TF-IDF baselines): lowercase, remove
        URLs, HTML tags, and all non-ASCII characters. This works well for
        English but destroys non-English text on purpose -- this shows why
        classical models fail across languages.
  \item \textbf{Light cleaning} (for LSTM and mBERT): remove only URLs, HTML
        tags, and extra whitespace. Unicode characters are kept so that Turkish,
        Spanish, and Italian text stays intact.
\end{itemize}

\begin{figure}[H]
  \centering
  \begin{subfigure}{0.48\textwidth}
    \includegraphics[width=\textwidth]{figures/label_distribution.png}
    \caption{Label distribution and toxic ratio per language.}
  \end{subfigure}
  \hfill
  \begin{subfigure}{0.48\textwidth}
    \includegraphics[width=\textwidth]{figures/text_length_by_toxicity.png}
    \caption{Comment length distribution by toxicity.}
  \end{subfigure}
  \caption{Exploratory data analysis of the Jigsaw dataset.}
  \label{fig:eda}
\end{figure}

% =========================================================================
\section{Literature Review}
% =========================================================================

\textbf{Multilingual transformers.}
Devlin et al.\ \cite{devlin2019} introduced BERT and its multilingual version
(mBERT), which is pretrained on 104 languages with a shared WordPiece
vocabulary. Conneau et al.\ \cite{conneau2020} proposed XLM-RoBERTa, which
improves cross-lingual transfer by using larger corpora and better language
sampling.

\textbf{Adapter-based fine-tuning.}
Houlsby et al.\ \cite{houlsby2019} proposed inserting small bottleneck adapter
layers into each transformer block. Only these adapter weights are trained while
the pretrained model is frozen. This reduces the number of trainable parameters
significantly. Pfeiffer et al.\ \cite{pfeiffer2020} showed that adapters work
well for cross-lingual tasks through AdapterHub. Our project follows this
approach.

\textbf{Toxic comment detection.}
The Jigsaw competition series on Kaggle showed that transformer models strongly
outperform classical methods for toxicity detection \cite{risch2020}. Top
solutions typically use transformer ensembles, translation-based data
augmentation, and pseudo-labeling.

% =========================================================================
\section{Methodology}
% =========================================================================

\subsection{Baseline Models}

\textbf{TF-IDF + Logistic Regression.}
We extract TF-IDF features (50,000 features, unigrams and bigrams, sublinear TF
scaling) and train a Logistic Regression classifier with L2 regularization.

\textbf{TF-IDF + Naive Bayes.}
Multinomial Naive Bayes with the same TF-IDF features. This is a probabilistic
baseline to compare against discriminative methods.

\textbf{TF-IDF + SVM.}
Linear SVM with TF-IDF features, wrapped with probability calibration (3-fold
cross-validation) so we can get probability outputs for AUC-ROC calculation.

\textbf{BiLSTM.}
A bidirectional LSTM with a learned embedding layer (200 dimensions), hidden
size of 128, max-pooling over time steps, and dropout of 0.3. The model uses
BCEWithLogitsLoss with a positive class weight of $\sim$9.4 to handle class
imbalance. We select the checkpoint with the best validation AUC-ROC.

\textbf{mBERT -- full fine-tuning.}
We fine-tune \texttt{bert-base-multilingual-cased} (178M parameters) on the
English training data with all parameters updated. This is the strongest
baseline. We use BCEWithLogitsLoss with class weighting and AdamW optimizer
with linear warmup schedule.

\subsection{Proposed Method: mBERT with Adapters}

Our proposed method inserts \textbf{Houlsby-style bottleneck adapters} into
every transformer block of mBERT. Each adapter consists of:

\begin{enumerate}  \item A down-projection: $\mathbb{R}^{768} \rightarrow \mathbb{R}^{64}$
  \item A GELU activation
  \item An up-projection: $\mathbb{R}^{64} \rightarrow \mathbb{R}^{768}$
  \item A residual connection (input is added back)
\end{enumerate}

Two adapters are placed in each transformer layer: one after the self-attention
output and one after the feed-forward output. The pretrained mBERT weights are
frozen. Only adapters, LayerNorm parameters, and the classification head are
trained. This gives us \textbf{2.4M trainable parameters} out of 180M total
(\textbf{1.34\%}).

The adapter weights are initialized near zero so that the adapter starts as an
approximate identity function and does not break the pretrained representations.

\begin{figure}[H]
  \centering
  \includegraphics[width=0.65\textwidth]{figures/param_inference_comparison.png}
  \caption{Trainable parameters and inference time comparison between
           mBERT-adapter and mBERT-full.}
  \label{fig:params}
\end{figure}

\subsection{Training Details}

All deep learning models use the same setup:
\begin{itemize}  \item Class-weighted BCEWithLogitsLoss (positive weight $\approx$ 9.4).
  \item Early stopping based on multilingual validation AUC-ROC.
  \item Decision threshold tuned on the multilingual validation set to
        maximize F1-score (instead of using the default 0.5).
\end{itemize}

% =========================================================================
\section{Results}
% =========================================================================

\subsection{English Test Results}

Table~\ref{tab:english} shows the results on the 20\% English holdout set.

\begin{table}[H]
\centering
\caption{Results on the English test split (20\% holdout).}
\label{tab:english}
\begin{tabular}{lccccc}
\toprule
\textbf{Model} & \textbf{AUC-ROC} & \textbf{F1} & \textbf{Accuracy} & \textbf{Precision} & \textbf{Recall} \\
\midrule
TF-IDF + Log.\ Reg.      & 0.9693 & 0.6993 & 0.9527 & 0.8907 & 0.5756 \\
TF-IDF + Naive Bayes      & 0.9434 & 0.6759 & 0.9445 & 0.7658 & 0.6049 \\
TF-IDF + SVM               & 0.9649 & 0.7299 & 0.9555 & 0.8712 & 0.6280 \\
BiLSTM                      & 0.9362 & 0.3587 & 0.6720 & 0.2206 & 0.9586 \\
mBERT-full                  & 0.9612 & 0.6918 & 0.9259 & 0.5746 & 0.8693 \\
\textbf{mBERT-adapter-64}   & \textbf{0.9732} & 0.6846 & 0.9200 & 0.5494 & 0.9081 \\
\bottomrule
\end{tabular}
\end{table}

On English, all models achieve high AUC-ROC ($>$0.93). The TF-IDF baselines
have the best accuracy because their high threshold makes few false positive
errors. The mBERT-adapter model achieves the \textbf{highest AUC-ROC (0.9732)},
showing that adapter fine-tuning does not hurt English performance.

The BiLSTM has very high recall (0.96) but very low precision (0.22), meaning
it flags most comments as toxic. This happens because the class-weighted loss
pushes the model too aggressively toward the minority class when the threshold
is low.

\subsection{Multilingual Results}

Table~\ref{tab:multi} shows the cross-lingual transfer results on the
multilingual validation set.

\begin{table}[H]
\centering
\caption{Results on the multilingual validation set (Turkish, Spanish, Italian).}
\label{tab:multi}
\begin{tabular}{lccccc}
\toprule
\textbf{Model} & \textbf{AUC-ROC} & \textbf{F1} & \textbf{Accuracy} & \textbf{Precision} & \textbf{Recall} \\
\midrule
TF-IDF + Log.\ Reg.      & 0.6266 & 0.0177 & 0.8472 & 0.7857 & 0.0089 \\
TF-IDF + Naive Bayes      & 0.6048 & 0.0222 & 0.8459 & 0.4516 & 0.0114 \\
TF-IDF + SVM               & 0.5816 & 0.0225 & 0.8478 & 0.8750 & 0.0114 \\
BiLSTM                      & 0.6005 & 0.2928 & 0.7023 & 0.2306 & 0.4008 \\
\textbf{mBERT-full}         & \textbf{0.8477} & \textbf{0.3964} & 0.8565 & 0.5610 & 0.3065 \\
mBERT-adapter-64            & 0.8261 & 0.3394 & 0.8506 & 0.5302 & 0.2496 \\
\bottomrule
\end{tabular}
\end{table}

The most important finding is the \textbf{huge gap between classical and
transformer models} on multilingual data. TF-IDF baselines drop to AUC-ROC
$\sim$0.58--0.63 and near-zero F1 because their ASCII-only features cannot
represent non-English text. The mBERT models achieve AUC-ROC $>$0.82, showing
successful cross-lingual transfer.

mBERT-full slightly outperforms mBERT-adapter on the multilingual validation
set (0.8477 vs.\ 0.8261 AUC-ROC). This is expected because full fine-tuning has
more capacity. However, the adapter model achieves \textbf{97.5\% of the full
model's AUC-ROC while training only 1.34\% of the parameters}.

\subsection{Per-Language Results}

\begin{figure}[H]
  \centering
  \includegraphics[width=0.75\textwidth]{figures/per_language_auc.png}
  \caption{Per-language AUC-ROC comparison across all models.}
  \label{fig:perlang}
\end{figure}

Table~\ref{tab:perlang} breaks down the mBERT results by language.

\begin{table}[H]
\centering
\caption{Per-language AUC-ROC for the mBERT models.}
\label{tab:perlang}
\begin{tabular}{lccc}
\toprule
\textbf{Model} & \textbf{Turkish} & \textbf{Spanish} & \textbf{Italian} \\
\midrule
mBERT-adapter-64 & 0.8900 & 0.8224 & 0.7719 \\
mBERT-full        & 0.9322 & 0.8374 & 0.7827 \\
\bottomrule
\end{tabular}
\end{table}

Both models work best on Turkish (AUC-ROC $>$0.89) and worst on Italian
($\sim$0.77--0.78). Turkish performance is strong possibly because Turkish has
more samples in the validation set (3,000 vs.\ 2,500 for others).

% =========================================================================
\section{Ablation Study}
% =========================================================================

\subsection{Adapter Size}

We used a bottleneck size of 64 in our main experiments. The adapter module has
a down-projection to 64 dimensions and an up-projection back to 768. With 24
adapters across 12 transformer layers, this gives 2.4M trainable parameters.
Larger bottleneck sizes (128, 256) would increase capacity but also increase
the parameter count.

\subsection{Full Fine-Tuning vs.\ Adapters}

\begin{table}[H]
\centering
\caption{Adapter vs.\ full fine-tuning comparison.}
\label{tab:ablation}
\begin{tabular}{lcccc}
\toprule
\textbf{Model} & \textbf{Trainable Params} & \textbf{\% of Total} & \textbf{EN AUC} & \textbf{Multi AUC} \\
\midrule
mBERT-adapter-64 & 2.4M  & 1.34\%  & 0.9732 & 0.8261 \\
mBERT-full        & 177.9M & 100\%  & 0.9612 & 0.8477 \\
\bottomrule
\end{tabular}
\end{table}

The adapter model trains \textbf{74$\times$ fewer parameters} than full
fine-tuning. On English it actually achieves a higher AUC-ROC (0.9732 vs.\
0.9612). On multilingual data, full fine-tuning is slightly better (0.8477 vs.\
0.8261). This shows that adapters are a very efficient trade-off: they lose only
a small amount of cross-lingual performance while being much cheaper to train.

\subsection{Text Cleaning Impact}

The TF-IDF baselines use aggressive ASCII-only cleaning, which removes all
non-English characters. This is the main reason they fail on multilingual data
(AUC-ROC drops from $\sim$0.96 to $\sim$0.60). The LSTM and mBERT models use
light cleaning that keeps Unicode characters, which is essential for
cross-lingual performance.

% =========================================================================
\section{Error Analysis}
% =========================================================================

\begin{figure}[H]
  \centering
  \begin{subfigure}{0.48\textwidth}
    \includegraphics[width=\textwidth]{figures/confusion_mBERT-adapter-64.png}
    \caption{mBERT-adapter confusion matrices.}
  \end{subfigure}
  \hfill
  \begin{subfigure}{0.48\textwidth}
    \includegraphics[width=\textwidth]{figures/confusion_SVM.png}
    \caption{TF-IDF + SVM confusion matrices.}
  \end{subfigure}
  \caption{Per-language confusion matrices for multilingual validation.}
  \label{fig:confusion}
\end{figure}

The confusion matrices in Figure~\ref{fig:confusion} show the main error
patterns:

\begin{itemize}  \item \textbf{TF-IDF models} predict almost everything as non-toxic on
        multilingual data. They have near-zero recall because they cannot
        recognize toxic words in Turkish, Spanish, or Italian.
  \item \textbf{mBERT-adapter} correctly identifies many toxic comments across
        languages but still misses some. The main error type is false negatives
        (toxic comments predicted as non-toxic).
  \item \textbf{BiLSTM} has the opposite problem: it flags too many comments as
        toxic (high recall but low precision) because it uses an aggressive
        threshold to compensate for its English-only vocabulary.
\end{itemize}

% =========================================================================
\section{Visualizations}
% =========================================================================

\begin{figure}[H]
  \centering
  \begin{subfigure}{0.48\textwidth}
    \includegraphics[width=\textwidth]{figures/roc_English_Test_20.png}
    \caption{ROC curves -- English test.}
  \end{subfigure}
  \hfill
  \begin{subfigure}{0.48\textwidth}
    \includegraphics[width=\textwidth]{figures/roc_Multilingual_Val.png}
    \caption{ROC curves -- Multilingual validation.}
  \end{subfigure}
  \caption{ROC curves comparing all models.}
  \label{fig:roc}
\end{figure}

\begin{figure}[H]
  \centering
  \includegraphics[width=0.85\textwidth]{figures/tsne_mbert.png}
  \caption{t-SNE visualization of mBERT [CLS] embeddings on the multilingual
           validation set. Left: colored by language. Right: colored by toxicity.
           Languages are mixed together (good cross-lingual alignment) and
           toxic/non-toxic clusters show some separation.}
  \label{fig:tsne}
\end{figure}

\begin{figure}[H]
  \centering
  \includegraphics[width=0.7\textwidth]{figures/tfidf_feature_importance.png}
  \caption{Top TF-IDF features for the toxic and non-toxic classes (from the
           Logistic Regression model). Toxic features include profanity and
           insults; non-toxic features are neutral Wikipedia-style words.}
  \label{fig:features}
\end{figure}

% =========================================================================
\section{GitHub and Reproducibility}
% =========================================================================

The complete source code is available at:
\begin{center}
  \url{https://github.com/yasntrk/multilingual-toxic-comment-classification}
\end{center}

The repository contains:
\begin{itemize}  \item \texttt{src/data\_utils.py} -- shared data loading and cleaning
  \item \texttt{src/metrics\_utils.py} -- shared evaluation metrics
  \item \texttt{src/eda.py} -- exploratory data analysis
  \item \texttt{src/train\_baselines.py} -- TF-IDF baselines
  \item \texttt{src/train\_lstm.py} -- BiLSTM training
  \item \texttt{src/train\_transformer.py} -- mBERT (full and adapter)
  \item \texttt{src/evaluate.py} -- result tables and visualizations
  \item \texttt{src/predict\_app.py} -- interactive demo
  \item \texttt{requirements.txt} -- all Python dependencies
  \item \texttt{results/} -- CSV/JSON files with all experiment results
  \item \texttt{figures/} -- all generated plots
\end{itemize}

To reproduce the experiments: install dependencies with
\texttt{pip install -r requirements.txt}, download the dataset from Kaggle, and
run the training scripts. The README file has full instructions.

% =========================================================================
\section{Conclusion}
% =========================================================================

In this project, we built a multilingual toxic comment classifier and compared
six different models. Our main findings are:

\begin{enumerate}  \item \textbf{Classical ML models fail across languages.} TF-IDF baselines
        achieve AUC-ROC $\sim$0.96 on English but drop to $\sim$0.60 on
        multilingual data because they cannot represent non-English text.
  \item \textbf{mBERT enables cross-lingual transfer.} Both mBERT models
        achieve AUC-ROC $>$0.82 on multilingual data without seeing any
        non-English training data.
  \item \textbf{Adapters are an efficient alternative.} The adapter model trains
        only 1.34\% of the parameters and reaches 97.5\% of full fine-tuning's
        multilingual AUC-ROC. On English, adapters even outperform full
        fine-tuning (0.9732 vs.\ 0.9612).
  \item \textbf{Class imbalance is a key challenge.} The 10\% toxic ratio
        causes low F1 scores across all models. Threshold tuning and class
        weighting help but do not fully solve the problem.
\end{enumerate}

For future work, we could try XLM-RoBERTa as the base model, experiment with
larger adapter sizes, use translation-based data augmentation, or try combining
adapters with full fine-tuning of the top layers.

% =========================================================================
\begin{thebibliography}{9}

\bibitem{jigsaw}
Jigsaw/Conversation AI,
``Jigsaw Multilingual Toxic Comment Classification,''
Kaggle, 2020.
\url{https://www.kaggle.com/c/jigsaw-multilingual-toxic-comment-classification}

\bibitem{devlin2019}
J.~Devlin, M.-W.~Chang, K.~Lee, and K.~Toutanova,
``BERT: Pre-training of deep bidirectional transformers for language understanding,''
\textit{NAACL-HLT}, 2019.

\bibitem{conneau2020}
A.~Conneau, K.~Khandelwal, N.~Goyal, et al.,
``Unsupervised cross-lingual representation learning at scale,''
\textit{ACL}, 2020.

\bibitem{houlsby2019}
N.~Houlsby, A.~Giurgiu, S.~Jastrzebski, et al.,
``Parameter-efficient transfer learning for NLP,''
\textit{ICML}, 2019.

\bibitem{pfeiffer2020}
J.~Pfeiffer, A.~Kamath, A.~R\"{u}ckl\'{e}, K.~Cho, and I.~Gurevych,
``AdapterHub: A framework for adapting transformers,''
\textit{EMNLP}, 2020.

\bibitem{risch2020}
J.~Risch, R.~Krestel,
``Toxic comment detection in online discussions,''
\textit{Deep Learning-Based Approaches for Sentiment Analysis}, Springer, 2020.

\end{thebibliography}

\end{document}
